\chapter{Ответы}
\label{ch:answers}

\openepigraph{%
Ставлю три звездочки. 
Я видал в детских книжках: 
когда человек делает прыжок к новой мысли, он ставит три звездочки...
}{Саша Чёрный, {Дневник Фокса Микки}
}

\newthought{Задачи рассеяны} по всей книге, а ответы на них собраны в этом разделе.

\begin{task}
    \label{answer:Tricolor-flags}
    \newthought{Меняя порядок и толщину} горизонтальных полосок 
    сербского флага (с.~\pageref{fig:Srpska}), можно получить флаг России 
    или флаг Крыма (рис.~\ref{fig:FlagCrimea}). 
\end{task}

\begin{marginfigure}[0cm]
{%
\setlength{\fboxsep}{0pt}%
\setlength{\fboxrule}{1pt}%
\fcolorbox{gray}{gray}{\includegraphics{./commons/Flag_of_Crimea.png}}%
}%
\caption{Флаг Крыма}
\label{fig:FlagCrimea}
\end{marginfigure}

\begin{task}
    \label{answer:Pobeda-beda}
    \newthought{Название корабля укоротилось} в мультфильме <<Приключения капитана 
    Врунгеля>>, поскольку юнга Лом срубил свежий лес для постройки яхты, поэтому корабль 
    и прирос к берегу. Пришлось рубить корни, первые две буквы не удержались и отвалились, 
    яхта <<Победа>> превратилась в <<Беду>>. 
    \small{Вопрос на с.~\pageref{fig:Pobeda-beda}.}
\end{task}

Задача: сколько ссылок можно закодировать, если ссылки содержит ровно 7 буквенно-цифровых символов,
     каждый символ может быть буквой английского алфавита 
     или цифрой от 0 до 9? Ответ на странице ... todo 

temp: Чем хороши глобальные переменные и какие у них недостатки
по сравнению с локальными переменными?

\begin{task}
    \label{answer:guess_numbers_task}
    \newthought{В данном алгоритме число $ a $ — натуральное}. 
    \small{Вопрос на с.~\pageref{question:text}}
\end{task}

\begin{task}
    \label{answer:naming}
    \newthought{Существуют следующие стили написания составных слов}: CamelCase, SnakeCase, KebabCase, PascalCase и другие. 
    \small{Упоминание на с.~\pageref{styles}}
\end{task}

\begin{task}
    \label{answer:global-vars-pros-cons}
    \newthought{Ответ такой} ... todo. 

    \small{Вопрос на с.~\pageref{fig:block:proc:swap:colors}.}
\end{task}

